\begin{abstract}
Object grounding requires a vision--language model to distinguish evidence of
an object's presence from contextual cues that merely make the object likely.
Reducing hallucinations by making a model more reluctant to answer positively
can obscure this distinction and increase omissions of genuinely visible
objects. We study whether a small intervention in attention-head outputs can
improve this tradeoff. We propose Evidence-Preserving Head Steering
(\method), which separates offline causal localization from supervised
intervention fitting. Localization contrasts activation-patching effects on
target-occluded images with effects on matched background controls. The
selected heads receive bounded scalar gains, fitted with balanced
object-existence supervision and a distribution-preservation objective while
the pretrained backbone remains frozen. Inference requires only the input
image and question, without donor activations or object annotations. We
specify an evaluation protocol that separates sensitivity to object presence
from sensitivity to surrounding context, alongside natural-image
hallucination, object recall, and capability measurements. The central
hypothesis is that control-adjusted head selection improves object grounding
beyond response-threshold calibration and equally sized arbitrary
interventions. This draft presents the implemented method and the protocol
for testing that hypothesis; empirical performance remains to be established.
\end{abstract}
